% !PRINTED 325
\weberpage{325}
%<<CONT
is called the discriminant of $f$. By this definition it is determined only up to a numerical factor, and up to now it can be computed only in the simplest cases.

The points of the curve $f$ whose coordinates satisfy the conditions (1) are called \gesperrt{singular points}.

One may also raise the question under what condition it is possible that infinitely many singular points occur on a curve. For this to happen, first of all the result of eliminating one of the unknowns, say $x_3$, from the two equations $f'(x_1) = 0$, $f'(x_2) = 0$ must vanish identically, i.\,e. the two functions $f'(x_1)$ and $f'(x_2)$, regarded as functions of $x_3$, must have a common divisor, and by Bd. I, \S.\,51 they must therefore also have a (non-constant) common divisor in the region of the forms $x_1, x_2, x_3$. This divisor must then in turn, as follows from the same considerations, have a common divisor with $f$, or, what is the same thing, with $f'(x_3)$, and it thus results that infinitely many singular points can be present only when the four forms $f, f'(x_1), f'(x_2), f'(x_3)$ have a common divisor.

Now let $u$ be such a common divisor, which we assume to be \gesperrt{irreducible}.

We set
\[
f(x_1, x_2, x_3) = uv
\]
\[
f'(x_1) = u\,\frac{\partial v}{\partial x_1} + v\,\frac{\partial u}{\partial x_1}
\]
\[
f'(x_2) = u\,\frac{\partial v}{\partial x_2} + v\,\frac{\partial u}{\partial x_2}
\]
\[
f'(x_3) = u\,\frac{\partial v}{\partial x_3} + v\,\frac{\partial u}{\partial x_3},
\]
and these equations show, since $u$ is irreducible and therefore
\[
u, \quad \frac{\partial u}{\partial x_1}, \quad \frac{\partial u}{\partial x_2}, \quad \frac{\partial u}{\partial x_3}
\]
can have no common divisor, that $v$ must be divisible by $u$; i.\,e. infinitely many singular points can occur only when $f(x_1, x_2, x_3)$ has a \gesperrt{quadratic divisor}. In this case we say that the curve $f$ contains a doubly counting component.
